% Compile with: pdflatex Exam_Study_Guide.tex
\documentclass[9pt,a4paper,twocolumn]{article}

\usepackage[utf8]{inputenc}
\usepackage[T5]{fontenc}
\usepackage[vietnamese]{babel}
\usepackage[margin=1.2cm,top=1.4cm,bottom=1.2cm,columnsep=0.5cm]{geometry}
\usepackage[table,dvipsnames]{xcolor}
\usepackage{tcolorbox}
\tcbuselibrary{breakable,skins}
\usepackage{booktabs}
\usepackage{enumitem}
\usepackage{listings}
\usepackage{microtype}
\usepackage{parskip}
\usepackage{titlesec}
\usepackage{array}
\usepackage{tabularx}
\usepackage{multicol}

% ---- Spacing ----
\setlength{\parskip}{2pt}
\setlength{\parindent}{0pt}
\setlist{nosep, leftmargin=1.2em, topsep=1pt, itemsep=0pt}
\titlespacing*{\section}{0pt}{5pt}{2pt}
\titlespacing*{\subsection}{0pt}{3pt}{1pt}
\titleformat{\section}{\normalfont\bfseries\small}{}{0em}{}
\titleformat{\subsection}{\normalfont\bfseries\footnotesize}{}{0em}{}

% ---- Colors ----
\definecolor{weekcolor}{RGB}{220,235,255}
\definecolor{weekframe}{RGB}{70,130,200}
\definecolor{trapcolor}{RGB}{255,235,210}
\definecolor{trapframe}{RGB}{210,120,40}
\definecolor{insightcolor}{RGB}{215,240,215}
\definecolor{insightframe}{RGB}{60,160,60}
\definecolor{deepcolor}{RGB}{235,220,250}
\definecolor{deepframe}{RGB}{130,70,180}
\definecolor{codebg}{RGB}{245,245,245}

% ---- tcolorbox styles ----
\tcbset{
  weekbox/.style={
    colback=weekcolor, colframe=weekframe, boxrule=0.7pt,
    fonttitle=\bfseries\small, coltitle=black,
    left=3pt, right=3pt, top=2pt, bottom=2pt,
    toptitle=2pt, bottomtitle=1pt, breakable
  },
  trapbox/.style={
    colback=trapcolor, colframe=trapframe, boxrule=0.5pt,
    left=3pt, right=3pt, top=2pt, bottom=2pt, breakable
  },
  insightbox/.style={
    colback=insightcolor, colframe=insightframe, boxrule=0.5pt,
    left=3pt, right=3pt, top=2pt, bottom=2pt, breakable
  },
  deepbox/.style={
    colback=deepcolor, colframe=deepframe, boxrule=0.7pt,
    fonttitle=\bfseries\small, coltitle=black,
    left=3pt, right=3pt, top=2pt, bottom=2pt,
    toptitle=2pt, bottomtitle=1pt, breakable
  },
}

% ---- Code listings ----
\lstset{
  basicstyle=\ttfamily\scriptsize,
  backgroundcolor=\color{codebg},
  frame=none, breaklines=true,
  columns=fullflexible, keepspaces=true,
  aboveskip=2pt, belowskip=2pt,
  xleftmargin=3pt, xrightmargin=3pt,
  inputencoding=utf8,
  extendedchars=true,
  literate=
    {á}{{\'a}}1 {à}{{\`a}}1 {ả}{{\h{a}}}1 {ã}{{\~a}}1 {ạ}{{\d{a}}}1
    {ă}{{\u{a}}}1 {ắ}{{\'\u{a}}}1 {ằ}{{\`\u{a}}}1 {ẳ}{{\h{\u{a}}}}1
    {ẵ}{{\~\u{a}}}1 {ặ}{{\d{\u{a}}}}1
    {â}{{\^a}}1 {ấ}{{\'{\^a}}}1 {ầ}{{\`{\^a}}}1 {ẩ}{{\h{\^a}}}1
    {ẫ}{{\~{\^a}}}1 {ậ}{{\d{\^a}}}1
    {đ}{{\dj}}1 {Đ}{{\DJ}}1
    {é}{{\'e}}1 {è}{{\`e}}1 {ẻ}{{\h{e}}}1 {ẽ}{{\~e}}1 {ẹ}{{\d{e}}}1
    {ê}{{\^e}}1 {ế}{{\'{\^e}}}1 {ề}{{\`{\^e}}}1 {ể}{{\h{\^e}}}1
    {ễ}{{\~{\^e}}}1 {ệ}{{\d{\^e}}}1
    {í}{{\'i}}1 {ì}{{\`i}}1 {ỉ}{{\h{i}}}1 {ĩ}{{\~i}}1 {ị}{{\d{i}}}1
    {ó}{{\'o}}1 {ò}{{\`o}}1 {ỏ}{{\h{o}}}1 {õ}{{\~o}}1 {ọ}{{\d{o}}}1
    {ô}{{\^o}}1 {ố}{{\'{\^o}}}1 {ồ}{{\`{\^o}}}1 {ổ}{{\h{\^o}}}1
    {ỗ}{{\~{\^o}}}1 {ộ}{{\d{\^o}}}1
    {ơ}{{\ohorn}}1 {ớ}{{\'{\ohorn}}}1 {ờ}{{\`{\ohorn}}}1
    {ở}{{\h{\ohorn}}}1 {ỡ}{{\~{\ohorn}}}1 {ợ}{{\d{\ohorn}}}1
    {ú}{{\'u}}1 {ù}{{\`u}}1 {ủ}{{\h{u}}}1 {ũ}{{\~u}}1 {ụ}{{\d{u}}}1
    {ư}{{\uhorn}}1 {ứ}{{\'{\uhorn}}}1 {ừ}{{\`{\uhorn}}}1
    {ử}{{\h{\uhorn}}}1 {ữ}{{\~{\uhorn}}}1 {ự}{{\d{\uhorn}}}1
    {ý}{{\'y}}1 {ỳ}{{\`y}}1 {ỷ}{{\h{y}}}1 {ỹ}{{\~y}}1 {ỵ}{{\d{y}}}1
    {→}{{$\rightarrow$}}1 {←}{{$\leftarrow$}}1
    {↓}{{$\downarrow$}}1 {↑}{{$\uparrow$}}1
    {↔}{{$\leftrightarrow$}}1
    {─}{{\textemdash}}1 {│}{{|}}1
    {┌}{{+}}1 {┐}{{+}}1 {└}{{+}}1 {┘}{{+}}1
    {├}{{+}}1 {┤}{{+}}1 {┬}{{+}}1 {┴}{{+}}1 {┼}{{+}}1
}

\newcommand{\code}[1]{\texttt{\small #1}}
\newcommand{\bold}[1]{\textbf{#1}}

% ===================================================================
\begin{document}
% ===================================================================

\begin{center}
  {\large\bfseries Ôn Thi --- Advanced IoT (7 tuần)}\\[1pt]
  {\footnotesize TS. Huỳnh Văn Đặng | UIT-VNUHCM | Đề mở, tự luận}
\end{center}
\vspace{-4pt}
\hrule
\vspace{4pt}

% -------------------------------------------------------------------
\section*{Mạch chung --- Một câu hỏi, bảy tuần}
% -------------------------------------------------------------------

\textit{``Làm sao xây IoT \textbf{thông minh, bảo mật, kết nối toàn cầu} khi device nhỏ,
pin ít, data nhạy cảm, mạng không phủ hết?''}

\begin{itemize}
  \item \textbf{W1}: Kiến trúc tổng thể? Digital Twin là đích đến.
  \item \textbf{W2}: AI chạy trên device nhỏ thế nào? --- TinyML
  \item \textbf{W3}: Học cộng tác không chia sẻ data? --- FL
  \item \textbf{W4}: Giao tiếp local không phụ thuộc cloud? --- Matter/Thread
  \item \textbf{W5}: Kết nối toàn cầu, vùng xa? --- 6G/NTN
  \item \textbf{W6}: Bảo vệ toàn stack? --- Security
  \item \textbf{W7}: Những gì chưa giải được? --- Open Research
\end{itemize}

% -------------------------------------------------------------------
\begin{tcolorbox}[weekbox, title={W1 --- Kiến trúc IoT \& Digital Twin}]
% -------------------------------------------------------------------

\textbf{Ba era IoT}

\begin{tabular}{@{}p{1.6cm}p{3.8cm}@{}}
\toprule
\scriptsize Era & \scriptsize Đặc điểm chính \\
\midrule
\scriptsize IoT 1.0 & \scriptsize Sensor→server, người nhìn dashboard \\
\scriptsize IoT 2.0 & \scriptsize Cloud analytics, scale lớn hơn \\
\scriptsize IoT 3.0 & \scriptsize Digital Twin, tự dự đoán, vòng lặp kín \\
\bottomrule
\end{tabular}

\vspace{3pt}
\textbf{Lưu ý quan trọng:} 1.0→2.0 paradigm KHÔNG đổi (vẫn human-in-loop).
Bước nhảy thật là \textbf{2.0→3.0}:

\begin{tabular}{@{}p{2cm}p{1.5cm}p{1.9cm}@{}}
\toprule
\scriptsize & \scriptsize 1.0 \& 2.0 & \scriptsize 3.0 (DT) \\
\midrule
\scriptsize Ai quyết định & \scriptsize Người & \scriptsize Hệ thống \\
\scriptsize Phản ứng & \scriptsize Sau khi xảy & \scriptsize Trước khi xảy \\
\scriptsize Cloud & \scriptsize Bắt buộc & \scriptsize Tuỳ chọn \\
\scriptsize Vòng lặp & \scriptsize Mở & \scriptsize Kín \\
\bottomrule
\end{tabular}

\vspace{3pt}
\textbf{Digital Twin = Đích, không phải mảng.} Mọi W2--W6 đều nuôi DT.
Khi thi: kết thúc câu nào cũng được bằng
\textit{``...tầng cần thiết để hiện thực hoá Digital Twin.''}

\vspace{3pt}
\textbf{AAS / DTDL / NGSI-LD} --- tại sao cần? DT không hiểu raw JSON.
Cần ngữ nghĩa: \texttt{t} là nhiệt độ hay thời gian? Đơn vị? Ngưỡng?

\begin{tabular}{@{}p{1.2cm}p{1.8cm}p{2.4cm}@{}}
\toprule
\scriptsize Chuẩn & \scriptsize Nguồn & \scriptsize Câu hỏi trung tâm \\
\midrule
\scriptsize AAS & \scriptsize Industry 4.0 & \scriptsize Vòng đời tài sản? \\
\scriptsize DTDL & \scriptsize Microsoft & \scriptsize Interface thiết bị? \\
\scriptsize NGSI-LD & \scriptsize ETSI SmartCity & \scriptsize Quan hệ với thế giới? \\
\bottomrule
\end{tabular}

\end{tcolorbox}

\begin{tcolorbox}[insightbox]
\scriptsize \textbf{Insight W1:} Mục tiêu IoT không phải thu thập data --- mà là
\textbf{hành động tự động từ data}. Digital Twin là não ảo giữa data và hành động.
\end{tcolorbox}

% -------------------------------------------------------------------
\begin{tcolorbox}[weekbox, title={W2 --- TinyML (Trí tuệ trên device nhỏ)}]
% -------------------------------------------------------------------

\textbf{Pain (cloud AI):} latency (100s ms), bandwidth, privacy, resilience.

\textbf{Giải pháp:} inference ON device (MCU vài trăm KB RAM).

\textbf{Bottleneck thật:} Peak \textbf{SRAM} (activations khi inference), không phải Flash (weights).
$\Rightarrow$ Quantization (FP32→INT8) + Pruning (cắt weights~$\approx$0).
$\Rightarrow$ \textbf{MCUNet:} co-design model + runtime từ đầu cho hardware cụ thể.

\end{tcolorbox}

\begin{tcolorbox}[trapbox]
\scriptsize \textbf{Trap holes W2:}
\begin{enumerate}
  \item TinyML $\neq$ nén model --- phải co-design từ đầu (MCUNet)
  \item \textbf{Chỉ inference, không training} --- training vẫn ở cloud
  \item Model tốt = \textbf{deployable}, không phải accurate nhất
  \item Không loại bỏ cloud, chỉ giảm phụ thuộc (hybrid)
  \item TinyML + FL là \textbf{cặp đôi}: TinyML dùng model; FL tạo model
\end{enumerate}
\end{tcolorbox}

% -------------------------------------------------------------------
\begin{tcolorbox}[weekbox, title={W3 --- Federated Learning (FL)}]
% -------------------------------------------------------------------

\textbf{Pain:} data phân tán, không thể gom (GDPR, bí mật KD, bandwidth).

\textbf{Cơ chế (SE view = Git):}
\begin{lstlisting}
fork W_global -> train local -> delta = W_local - W_global
-> push delta -> server: W_new = W_global + avg(all deltas)
\end{lstlisting}

\textbf{FL trong IoT:} MCU không train được (cần 3--10$\times$ RAM so với inference).
$\Rightarrow$ \textbf{FL client = Edge Gateway}, không phải sensor.
Sensor gửi raw data lên Edge. Edge train và gửi delta lên Cloud.

\textbf{Privacy stack} (cần cả 3 tầng):
\begin{itemize}
  \item Tầng 1 --- FL: data không rời device
  \item Tầng 2 --- Secure Aggregation: server chỉ thấy \textbf{tổng}, không thấy delta riêng
  \item Tầng 3 --- Differential Privacy: thêm nhiễu, chặn gradient inversion
\end{itemize}

\end{tcolorbox}

\begin{tcolorbox}[trapbox]
\scriptsize \textbf{Trap holes W3:}
\begin{enumerate}
  \item FL $\neq$ privacy đảm bảo (cần SecAgg + DP)
  \item \textbf{SecAgg $\neq$ DP:} SecAgg chặn server-curious; DP chặn reconstruction attack
  \item \textbf{Non-IID} = vấn đề cốt lõi chưa giải: FedAvg non-IID $\approx$78\% vs centralized 95\%
  \item FL $\neq$ Distributed Training (privacy vs performance --- khác mục đích)
\end{enumerate}
\end{tcolorbox}

\begin{tcolorbox}[insightbox]
\scriptsize \textbf{Insight W3:} Chia sẻ tri thức (model), không chia sẻ data.
FL là câu trả lời kiến trúc cho bài toán privacy phân tán.
\end{tcolorbox}

% -------------------------------------------------------------------
\begin{tcolorbox}[weekbox, title={W4 --- Matter, Thread, MQTT, CoAP}]
% -------------------------------------------------------------------

\textbf{``Tháp Babel''} (Kinh Thánh): mỗi người nói tiếng khác $\rightarrow$ không ai hiểu ai $\rightarrow$ công trình sụp.
IoT trước 2020: Zigbee/Z-Wave/HAP/SmartThings = chiến lược \textbf{lock-in}.

\textbf{Matter:} IP-first, local-first, Multi-Admin.
Apple+Google+Amazon+Samsung đồng ý chuẩn chung tầng application.

\textbf{BLE trong Matter --- 2 lần xuất hiện:}
\begin{enumerate}
  \item Phần \textbf{vấn đề}: BLE truyền thống point-to-point, mesh đến muộn
  \item Phần \textbf{giải pháp}: BLE = bootstrap commissioning (sau đó tắt hẳn)
\end{enumerate}

\textbf{Commissioning flow:} BLE adv $\rightarrow$ QR scan (discriminator+passcode)
$\rightarrow$ PASE handshake (mã hoá BLE) $\rightarrow$ Device Attestation (cert nhà SX)
$\rightarrow$ gửi WiFi/Thread creds $\rightarrow$ BLE ngắt $\rightarrow$ Fabric assignment.

\textbf{Tại sao BLE?} Chicken-and-egg: cần BLE để truyền WiFi password vào device chưa có WiFi.

\textbf{Thread:} IPv6 mesh tự vá. Roles: Leader/Router/End Device.
Border Router kết nối mesh ra WiFi/Internet.

\vspace{3pt}
\begin{tabular}{@{}p{1.5cm}p{1.8cm}p{2.1cm}@{}}
\toprule
\scriptsize & \scriptsize MQTT & \scriptsize CoAP \\
\midrule
\scriptsize Mô hình & \scriptsize Pub/Sub (broker) & \scriptsize Req/Res (trực tiếp) \\
\scriptsize Transport & \scriptsize TCP & \scriptsize UDP \\
\scriptsize Security & \scriptsize TLS & \scriptsize DTLS + OSCORE \\
\scriptsize Tốt cho & \scriptsize Nhiều device→cloud & \scriptsize Device↔device \\
\bottomrule
\end{tabular}

\vspace{2pt}
Thực tế: \textbf{CoAP} device→edge, \textbf{MQTT} edge→cloud. Không phải đối thủ.
OSCORE \textbf{chỉ chạy với CoAP} (end-to-end qua proxy).

\end{tcolorbox}

\begin{tcolorbox}[insightbox]
\scriptsize \textbf{Insight W4:} Interoperability phải thiết kế upfront qua chuẩn.
``Local-first'' không phải bước lùi --- cloud là extension, không phải dependency.
\end{tcolorbox}

% -------------------------------------------------------------------
\begin{tcolorbox}[weekbox, title={W5 --- 6G \& Non-Terrestrial Networks (NTN)}]
% -------------------------------------------------------------------

\textbf{Tại sao W5 trong khoá IoT?} DT cần always-on connectivity.
W4 phủ local; W5 phủ sensor giữa đại dương/vùng núi/sa mạc.

\begin{tabular}{@{}p{0.8cm}p{2.0cm}p{2.6cm}@{}}
\toprule
\scriptsize & \scriptsize Bài toán giải & \scriptsize Cho ai \\
\midrule
\scriptsize 3G & \scriptsize Internet trên mobile & \scriptsize Người \\
\scriptsize 4G & \scriptsize Mobile broadband & \scriptsize Người \\
\scriptsize 5G & \scriptsize 3 nhu cầu khác nhau & \scriptsize Người + máy \\
\scriptsize 6G & \scriptsize Phủ sóng toàn cầu + AI & \scriptsize Máy móc chính \\
\bottomrule
\end{tabular}

\vspace{3pt}
\textbf{5G = 3 slices:} eMBB (video, high-BW), mMTC (IoT sensors, pin năm),
URLLC ($<$1ms, robot, xe tự lái).

\textbf{6G = 4 tầng NTN:}
\begin{itemize}
  \item Vệ tinh LEO (500--2000km, latency $\sim$20ms, toàn cầu)
  \item \textbf{HAPS} ($\sim$20km tầng bình lưu, quasi-stationary, vùng rộng)
  \item UAV (linh hoạt, on-demand base station)
  \item Mặt đất (dày đặc, đô thị)
\end{itemize}

\textbf{UAV -- 2 vai trò:} (1) user = kết nối như điện thoại nhưng có interference issues;
(2) \textbf{base station} = bay lên phủ sóng IoT bên dưới, backhaul qua satellite.

\textbf{HAPS:} trên mây + máy bay; dưới vệ tinh; phủ hàng trăm km; quasi-stationary.
(Ví dụ: Google Loon, Airbus Zephyr).

\textbf{3GPP:} Rel-17 = milestone tích hợp satellite; Rel-19 = cầu nối 6G.

\end{tcolorbox}

\begin{tcolorbox}[insightbox]
\scriptsize \textbf{Insight W5:} IoT toàn cầu không thể chỉ dựa mặt đất.
6G = AI-native + phủ sóng 3D (đất + trời + vũ trụ). NTN là kiến trúc chính, không phải add-on.
\end{tcolorbox}

% -------------------------------------------------------------------
\begin{tcolorbox}[weekbox, title={W6 --- IoT Security (Toàn stack)}]
% -------------------------------------------------------------------

\textbf{IT vs IoT Security:}

\begin{tabular}{@{}p{1.6cm}p{1.8cm}p{2.0cm}@{}}
\toprule
\scriptsize Yếu tố & \scriptsize IT & \scriptsize IoT \\
\midrule
\scriptsize Vòng đời & \scriptsize 3--5 năm & \scriptsize 10--20 năm, không vá \\
\scriptsize Vị trí & \scriptsize Data center & \scriptsize Ngoài trời tự do \\
\scriptsize Hậu quả & \scriptsize Mất data & \scriptsize Vật lý: tính mạng \\
\scriptsize Tài nguyên & \scriptsize TLS/RSA ok & \scriptsize MCU quá yếu \\
\bottomrule
\end{tabular}

\vspace{3pt}
\textbf{Attack vectors (lý do phải bảo vệ từng tầng):}
\begin{itemize}
  \item HW: firmware dump (JTAG), device clone, side-channel attack
  \item OTA: fake firmware push $\rightarrow$ kiểm soát fleet
  \item Network: MITM, replay attack (bắt ``mở cửa'' gửi lại), spoofing
  \item App: data poisoning (DT ra lệnh sai), command injection
  \item Scale: \textbf{Mirai 2016} = 600K cameras (default password) $\rightarrow$ DDoS Twitter/Netflix
  \item Cascade: device $\rightarrow$ gateway $\rightarrow$ cloud $\rightarrow$ toàn fleet
\end{itemize}

\textbf{Security stack (phải có đủ 4 tầng):}

\begin{itemize}
  \item \textbf{PUF}: vân tay vật lý chip (manufacturing variation), key sinh từ vật lý, không lưu bộ nhớ $\rightarrow$ không đọc được, không clone được
  \item \textbf{Secure Boot}: ROM $\rightarrow$ Bootloader $\rightarrow$ Firmware $\rightarrow$ App, mỗi bước verify chữ ký $\rightarrow$ firmware giả không chạy được
  \item \textbf{TEE/TrustZone}: Normal World (OS, có thể bị hack) $\Vert$ Secure World (keys, crypto). Key không bao giờ ra ngoài TEE.
  \item \textbf{Thread AES}: AES-128 link layer, hop-by-hop trong mesh, chống radio scanner
  \item \textbf{DTLS}: TLS cho UDP (CoAP dùng UDP). Thêm: seq number, retransmit, Connection ID
  \item \textbf{OSCORE}: mã hoá \textbf{message} (object), không phải channel. Sống qua proxy. DTLS/TLS=VPN tunnel; OSCORE=Signal E2E.
\end{itemize}

\begin{lstlisting}
[HW]  PUF -> Secure Boot -> TEE
[L2]  Thread AES (hop-by-hop)
[L4]  DTLS (per-segment)
[App] OSCORE (true end-to-end, qua proxy)
\end{lstlisting}

\textbf{STRIDE:} \textbf{S}poofing / \textbf{T}ampering / \textbf{R}epudiation /
\textbf{I}nfo Disclosure / \textbf{D}oS / \textbf{E}levation of Privilege

\end{tcolorbox}

\begin{tcolorbox}[insightbox]
\scriptsize \textbf{Insight W6:} Security không phải add-on --- phải là kiến trúc từ hardware.
Một tầng thủng $\Rightarrow$ cascade compromise. IoT security $\neq$ IT security thu nhỏ.
\end{tcolorbox}

% -------------------------------------------------------------------
\begin{tcolorbox}[weekbox, title={W7 --- Open Research (Biên giới tri thức)}]
% -------------------------------------------------------------------

W7 không dạy giải pháp --- dạy biết \textbf{ranh giới} của tri thức hiện tại.

\begin{itemize}
  \item \textbf{DT}: 1M sensors = 10M updates/s, physics model accuracy vs speed
  \item \textbf{TinyML}: fleet model management, model drift, MLOps cho IoT scale chưa có
  \item \textbf{FL}: non-IID gap (FedAvg 78\% vs centralized 95\%), heterogeneous devices, comm overhead
  \item \textbf{Matter}: legacy gap (tỉ Zigbee/Z-Wave devices), industrial IoT ngoài scope
  \item \textbf{6G}: THz (range vài chục m, bị air hấp thụ), RIS (cần exact position), ISAC (signal separation cực khó)
  \item \textbf{Security}: Zero Trust cho MCU --- continuous auth = pin cạn, history = RAM hết, policy enforcement = cần online liên tục
\end{itemize}

\end{tcolorbox}

\begin{tcolorbox}[insightbox]
\scriptsize \textbf{Insight W7:} IoT chưa phải ngành solved. Mỗi tuần học = giải pháp tốt nhất
hiện tại. W7 = biên giới, không phải cuối đường.
Câu tự luận hoàn chỉnh = giải pháp + trade-off + open challenge.
\end{tcolorbox}

% -------------------------------------------------------------------
\section*{Bản đồ kết nối 7 tuần}
% -------------------------------------------------------------------

\begin{lstlisting}
[W1:DT] <-data- [W4:Matter/Thread] <-extend- [W5:6G/NTN]
   |                   |
  AI              connectivity
   v                   v
[W2:TinyML] --FL-> [W3:FL]
   |                   |
   +---- both need -->[W6:Security]
                         |
                    all have gaps
                         v
                   [W7:Open Research]
\end{lstlisting}

% -------------------------------------------------------------------
\begin{tcolorbox}[deepbox, title={Framework Trả Lời Tự Luận (5 bước)}]
% -------------------------------------------------------------------
\begin{enumerate}
  \item \textbf{VẤN ĐỀ:} Pain point trước X?
  \item \textbf{GIẢI PHÁP:} Cơ chế X là gì?
  \item \textbf{ĐÁNH ĐỔI:} Trade-off? Khi nào không phù hợp?
  \item \textbf{KẾT NỐI:} Liên quan tuần nào?
  \item \textbf{OPEN:} Chưa giải được gì? (W7)
\end{enumerate}

\textbf{Ví dụ --- ``Tại sao cần FL?''}
\begin{enumerate}
  \item Data phân tán, không thể gom (GDPR/bandwidth/privacy)
  \item Edge Gateway train local $\rightarrow$ gửi delta $\rightarrow$ FedAvg aggregate
  \item Non-IID gap; gradient leakage (cần DP+SecAgg); MCU không train được
  \item TinyML (W2) dùng model FL tạo ra; W4/W5 truyền updates; W6 DP+SecAgg
  \item Non-IID convergence; fleet management; comm overhead (W7)
\end{enumerate}

\end{tcolorbox}

% ===================================================================
% DEEP DIVES
% ===================================================================
\vspace{4pt}
\hrule
\vspace{2pt}
\begin{center}{\small\bfseries Phân Tích Chuyên Sâu}\end{center}
\vspace{2pt}
\hrule
\vspace{4pt}

% -------------------------------------------------------------------
\begin{tcolorbox}[deepbox, title={Kiến trúc 4 tầng IoT --- Chi tiết}]
% -------------------------------------------------------------------

\begin{tabular}{@{}p{1.6cm}p{3.8cm}@{}}
\toprule
\scriptsize Tầng & \scriptsize Vai trò \\
\midrule
\scriptsize Device/Sensor & \scriptsize Thu thập data, actuator. RAM vài trăm KB, pin, CPU yếu. Protocols: Zigbee/BLE/Thread/LoRa \\
\scriptsize Edge/Gateway & \scriptsize Bộ não cục bộ. Tổng hợp, dịch protocol (Zigbee→IP), ra quyết định real-time. FL client thực sự. \\
\scriptsize Cloud & \scriptsize Lịch sử dài hạn, train ML, host Digital Twin, fleet management \\
\scriptsize User App & \scriptsize Dashboard/app, điểm can thiệp của con người \\
\bottomrule
\end{tabular}

\vspace{3pt}
\textbf{Hai luồng data khác nhau:}
\begin{lstlisting}
Analytics: Device->Edge->Cloud->App
Real-time: Device->Edge->Device (cloud khong tham gia)
\end{lstlisting}
Edge là nơi hai luồng tách ra.

\end{tcolorbox}

% -------------------------------------------------------------------
\begin{tcolorbox}[deepbox, title={Edge vs Cloud --- Framework ra quyết định}]
% -------------------------------------------------------------------

Phân tích theo \textbf{3 chiều} cho mọi use case:

\begin{tabular}{@{}p{1.6cm}p{1.8cm}p{2.0cm}@{}}
\toprule
\scriptsize Chiều & \scriptsize Edge & \scriptsize Cloud \\
\midrule
\scriptsize \textbf{Latency} & \scriptsize $<$100ms (safety) & \scriptsize Batch/daily \\
\scriptsize \textbf{Bandwidth} & \scriptsize Video/audio/lidar & \scriptsize Sensor nhỏ \\
\scriptsize \textbf{Computation} & \scriptsize Threshold, TinyML & \scriptsize Train model \\
\bottomrule
\end{tabular}

\vspace{3pt}
\textbf{Ví dụ: Nông nghiệp IoT}

\begin{tabular}{@{}p{2.4cm}p{0.8cm}p{0.8cm}p{0.5cm}p{0.6cm}@{}}
\toprule
\scriptsize Tác vụ & \scriptsize Lat & \scriptsize BW & \scriptsize CPU & \scriptsize $\Rightarrow$ \\
\midrule
\scriptsize Phát hiện sâu→phun & \scriptsize Gấp & \scriptsize Nặng & \scriptsize Vừa & \scriptsize \textbf{Edge} \\
\scriptsize Tưới theo lịch & \scriptsize Không & \scriptsize Nhỏ & \scriptsize Nhẹ & \scriptsize Cloud \\
\scriptsize Train model mới & \scriptsize Không & \scriptsize Lớn & \scriptsize Nặng & \scriptsize Cloud \\
\scriptsize Dự báo mùa vụ & \scriptsize Không & \scriptsize Lớn & \scriptsize Rất nặng & \scriptsize Cloud \\
\bottomrule
\end{tabular}

\vspace{3pt}
$\Rightarrow$ Kết luận luôn là \textbf{hybrid}: edge = real-time + data nặng; cloud = analytics + training.

\end{tcolorbox}

% -------------------------------------------------------------------
\begin{tcolorbox}[deepbox, title={Chọn Protocol Truyền Thông}]
% -------------------------------------------------------------------

\begin{lstlisting}
small + rare + remote   -> LoRaWAN / NB-IoT
small + local + battery -> Thread + CoAP
many devices -> cloud   -> MQTT
device <-> device       -> CoAP
video + local           -> WiFi 6
video + wide-area       -> 5G eMBB
<1ms critical control   -> 5G URLLC
\end{lstlisting}

\textbf{Tại sao Thread/CoAP/MQTT không dùng cho video?}
Thread $\approx$250Kbps --- một frame HD đã vượt quá.
MQTT/CoAP phụ thuộc bandwidth của mạng bên dưới, không phải protocol.

5G phân 3 slices vì \textbf{không thể tối ưu mMTC + eMBB + URLLC cùng lúc}.

\end{tcolorbox}

% -------------------------------------------------------------------
\begin{tcolorbox}[deepbox, title={Trade-off: Latency vs Accuracy}]
% -------------------------------------------------------------------

\textbf{Bản chất:}
\begin{lstlisting}
Model lon -> cloud -> latency cao -> accuracy cao
Model nho -> edge -> latency thap -> accuracy thap hon
\end{lstlisting}

\textbf{Ưu tiên latency khi:} safety-critical (sương giá, rò khí, dừng khẩn cấp).
Model edge 88\% + 50ms $>$ model cloud 95\% + 2s khi mái che đóng muộn = mùa màng thiệt.

\textbf{Ưu tiên accuracy khi:} chẩn đoán bệnh, dự báo bảo trì --- sai một lần hậu quả lớn hơn chậm.

\textbf{Cải thiện accuracy không hi sinh latency:}
\begin{itemize}
  \item \textbf{FL (W3):} nhiều edge train $\rightarrow$ model toàn cục tốt hơn theo thời gian; inference vẫn ở edge
  \item \textbf{MCUNet (W2):} co-design cho hardware cụ thể $\rightarrow$ accuracy tốt hơn trong cùng RAM
  \item \textbf{Hierarchical:} edge xử lý cases rõ ràng (high confidence); cloud xử lý edge cases
\end{itemize}

\end{tcolorbox}

% -------------------------------------------------------------------
\begin{tcolorbox}[deepbox, title={3 Điểm Yếu Kiến Trúc 4 Tầng (Security)}]
% -------------------------------------------------------------------

\textbf{1. Device tầng dưới: không bảo vệ được vật lý}

Server: data center có khoá. IoT device: ngoài đường, trong rừng, trên tàu.
Kẻ tấn công cầm tay $\rightarrow$ firmware dump, side-channel, device clone.
Vòng đời 20 năm $\rightarrow$ lỗ hổng sau không thể vá.

\textit{Attack surface này không tồn tại trong IT.}

\textbf{2. Gateway = single point of compromise}

Gateway biết credentials của \textbf{tất cả device bên dưới VÀ cloud bên trên}.
\begin{lstlisting}
Gateway bi hack -> pivot ca 2 huong:
    tat ca device + toan bo cloud
\end{lstlisting}
Single point of failure không có trong IT phân tán.

\textbf{3. Cascade attack path dọc}

Các tầng kết nối trực tiếp, không đủ isolation:
\begin{lstlisting}
Device (yeu nhat) -> Gateway -> Cloud -> fleet
\end{lstlisting}
\textbf{Mirai 2016:} camera default password $\rightarrow$ 600K botnet $\rightarrow$ DDoS Twitter/Netflix.

$\Rightarrow$ Phải bảo vệ từng tầng độc lập + mã hoá xuyên suốt (DTLS+OSCORE).

\end{tcolorbox}

% -------------------------------------------------------------------
\begin{tcolorbox}[deepbox, title={IoT Security vs IT Security --- Không thể dùng lại nguyên si}]
% -------------------------------------------------------------------

\begin{tabular}{@{}p{2.8cm}p{2.6cm}@{}}
\toprule
\scriptsize Assumption IT security & \scriptsize Thực tế IoT \\
\midrule
\scriptsize Đủ mạnh chạy TLS/RSA & \scriptsize MCU vài KB RAM $\rightarrow$ dùng DTLS, ECC, Ascon \\
\scriptsize Được vá thường xuyên & \scriptsize Vòng đời 15--20 năm, không ai vá $\rightarrow$ phải đúng từ đầu \\
\scriptsize Môi trường kiểm soát & \scriptsize Device ngoài trời $\rightarrow$ cần PUF, TEE \\
\scriptsize Vài nghìn thiết bị & \scriptsize Triệu thiết bị $\rightarrow$ key mgmt scale chưa giải \\
\scriptsize Hậu quả: mất data & \scriptsize Hậu quả: \textbf{máy nổ, tính mạng} \\
\scriptsize Protocol chuẩn hoá & \scriptsize Chục protocols, mỗi cái security khác \\
\bottomrule
\end{tabular}

\vspace{3pt}
\textbf{Điểm mấu chốt:} Mỗi dòng trong bảng là lý do một tầng trong security stack tồn tại.
PUF $\leftarrow$ device bị cầm tay. Secure Boot $\leftarrow$ không ai vá firmware.
DTLS $\leftarrow$ TLS quá nặng. OSCORE $\leftarrow$ proxy terminate DTLS.

\begin{tcolorbox}[insightbox]
\scriptsize IoT security $\neq$ IT security thu nhỏ --- đây là bài toán thiết kế \textbf{khác hoàn toàn},
đòi hỏi security-by-design từ phần cứng lên.
\end{tcolorbox}

\end{tcolorbox}

\end{document}
